% ===== Preâmbulo compartilhado (Tectonic / XeTeX) — documentos da FTL094 =====
\documentclass[11pt,a4paper]{article}
\usepackage{fontspec}
\usepackage[brazil]{babel}
\usepackage{geometry}
\geometry{a4paper,margin=2.5cm}
\usepackage{amsmath,amssymb}
\usepackage{graphicx}
\usepackage{booktabs}
\usepackage{array}
\usepackage{longtable}
\usepackage{tabularx}
\usepackage{float}
\usepackage{enumitem}
\usepackage{xcolor}
\usepackage{tikz}
\usetikzlibrary{arrows.meta,positioning,automata,shapes.geometric,fit,backgrounds,calc}
\usepackage{pgfplots}
\pgfplotsset{compat=1.18}
\usepackage{listings}
\usepackage{caption}
\usepackage{fancyhdr}
\usepackage[hidelinks]{hyperref}

% ----- cores -----
\definecolor{ufamblue}{RGB}{0,51,102}
\definecolor{codegreen}{rgb}{0,0.5,0}
\definecolor{codeblue}{rgb}{0,0,0.7}
\definecolor{codered}{rgb}{0.6,0,0}
\definecolor{lgray}{rgb}{0.95,0.95,0.95}

% ----- listagens C -----
\lstdefinestyle{c}{language=C,basicstyle=\ttfamily\footnotesize,
  keywordstyle=\color{codeblue}\bfseries,commentstyle=\color{codegreen}\itshape,
  stringstyle=\color{codered},numbers=left,numberstyle=\tiny\color{gray},
  frame=single,breaklines=true,tabsize=2,showstringspaces=false,columns=fullflexible}

% ----- constantes do projeto -----
\newcommand{\projnome}{Sistema de Navegação Autônoma para o Robô Móvel Khepera~IV}
\newcommand{\disciplina}{FTL094 --- Engenharia de Software de Tempo Real}
\newcommand{\periodo}{Período Letivo 2026/1}
\newcommand{\versaodoc}{1.0}

% ----- cabeçalho/rodapé -----
\pagestyle{fancy}\fancyhf{}
\rhead{\small \tipodoc}
\lhead{\small Khepera~IV --- FTL094}
\cfoot{\thepage}
\renewcommand{\headrulewidth}{0.4pt}

% ----- coluna X centralizada p/ tabularx -----
\newcolumntype{C}{>{\centering\arraybackslash}X}

% ----- capa padronizada -----
\newcommand{\fazcapa}{%
\begin{titlepage}\centering
{\large\bfseries UNIVERSIDADE FEDERAL DO AMAZONAS (UFAM)}\\[3pt]
{\normalsize Centro de P\&D em Tecnologia Eletrônica e da Informação --- CETELI}\\[2pt]
{\normalsize Programa de Pós-Graduação em Engenharia Elétrica --- PPGEE}\\[26pt]
{\large \disciplina}\\[2pt]
{\normalsize \periodo}\\[64pt]
\rule{\textwidth}{1.2pt}\\[12pt]
{\Huge\bfseries\color{ufamblue} \tipodoc}\\[14pt]
{\Large \projnome}\\[10pt]
\rule{\textwidth}{1.2pt}\\[48pt]
{\large\bfseries Equipe}\\[8pt]
{\normalsize Gabriel Henrique de Souza Nazaré}\\[3pt]
{\normalsize Luan Alexandre Ramos Barreto}\\[3pt]
{\normalsize Matheus Rocha Canto}\\[3pt]
{\normalsize Matheus Serrão Uchôa}\\[3pt]
{\normalsize Ricardo Mendonça Braz}\\[3pt]
{\normalsize Yan Fernandes da Silva}\\[40pt]
{\normalsize Docente responsável: Prof. Dr.~[nome do docente]}\\
\vfill
{\small Documento de tipo \emph{\tipodoc} --- Versão \versaodoc}\\[2pt]
{\small Manaus --- AM, \datadoc}
\end{titlepage}}

% ----- tabela de controle de versões -----
\newcommand{\controleversoes}{%
\section*{Controle de versões}
\begin{tabularx}{\textwidth}{@{}l l l X@{}}
\toprule
\textbf{Versão} & \textbf{Data} & \textbf{Autor} & \textbf{Descrição} \\
\midrule
1.0 & \datadoc & Equipe FTL094 & Emissão inicial do documento. \\
\bottomrule
\end{tabularx}
\bigskip}

\newcommand{\tipodoc}{Protocolo de Testes}
\newcommand{\datadoc}{30 de junho de 2026}

\begin{document}
\fazcapa
\controleversoes
\tableofcontents
\newpage

% =====================================================================
\section{Introdução}
% =====================================================================
Este documento apresenta os \textbf{resultados obtidos} na execução dos casos definidos no \emph{Plano de Testes} do \emph{\projnome}. Registra o veredito de cada caso (PASSOU/FALHOU), as evidências coletadas e os defeitos identificados e tratados durante a campanha de testes.

% =====================================================================
\section{Resumo da execução}
% =====================================================================
\begin{table}[H]\centering
\caption{Sumário da campanha de testes.}
\begin{tabularx}{\textwidth}{@{}l X@{}}
\toprule
\textbf{Item} & \textbf{Valor} \\
\midrule
Software testado & Controlador \texttt{goto\_avoid} (versão final, FSM Navegar/Girar/Contornar). \\
Ambiente & Webots R2025a, modo \emph{headless} (\texttt{--batch --mode=fast --no-rendering}). \\
Cenário & Arena $3\times3$~m; 5 obstáculos; início $(-1,-1)$; alvo $(1,1)$. \\
Período de execução & Junho/2026. \\
Casos executados & 6 (CT-01 a CT-06). \\
Veredito global & \textbf{APROVADO} (6/6 casos com resultado PASSOU). \\
\bottomrule
\end{tabularx}
\end{table}

% =====================================================================
\section{Resultados por caso de teste}
% =====================================================================
\begin{table}[H]\centering
\caption{Resultados da execução dos casos de teste.}
\renewcommand{\arraystretch}{1.25}
\begin{tabularx}{\textwidth}{@{}l X l X@{}}
\toprule
\textbf{ID} & \textbf{Descrição} & \textbf{Result.} & \textbf{Evidência} \\
\midrule
CT-01 & Navegação em caminho livre. & PASSOU & Segmento livre da missão ($t=0$--$2$\,s): avanço direto, distância $2{,}83\to2{,}40$~m. \\
CT-02 & Desvio de obstáculo frontal isolado. & PASSOU & Contorno da 1ª caixa ($t\approx3$\,s), estado \textsc{Contornar}. \\
CT-03 & Missão no \emph{slalom} de 5 obstáculos. & PASSOU & Chegada em $(0{,}94;\,1{,}03)$ em $19{,}3$\,s, \textbf{0 colisões} (Fig.~\ref{fig:traj}, \ref{fig:dist}). \\
CT-04 & Escape de mínimo local. & PASSOU$^\dagger$ & Após correção (DEF-01): FSM contorna; posição não estaciona. \\
CT-05 & Período do laço de controle. & PASSOU & Laço periódico de 16~ms; passos regulares na telemetria, sem perda de passo. \\
CT-06 & Telemetria e calibração de sensores. & PASSOU & \texttt{run.log} gerado; base de IR medida $\approx$110 (Tab.~\ref{tab:calib}). \\
\bottomrule
\end{tabularx}
\end{table}
\footnotesize $^\dagger$ Caso reprovado na primeira execução (defeito DEF-01) e \emph{reaprovado} após correção. \normalsize

% =====================================================================
\section{Evidências}
% =====================================================================
\subsection{CT-03 --- trajetória e convergência}
A Figura~\ref{fig:traj} mostra a trajetória livre de colisões sobre o mapa de obstáculos; a Figura~\ref{fig:dist}, a redução (quase) monotônica da distância ao alvo.

\begin{figure}[H]\centering
\begin{tikzpicture}
\begin{axis}[axis equal image,width=0.62\textwidth,xmin=-1.6,xmax=1.6,ymin=-1.6,ymax=1.6,
  xlabel={$x$ (m)},ylabel={$y$ (m)},grid=both,grid style={gray!20},title={CT-03: trajetória executada}]
  \draw[thick,gray](axis cs:-1.5,-1.5)rectangle(axis cs:1.5,1.5);
  \fill[brown!55](axis cs:-0.6,-0.6)rectangle(axis cs:-0.4,-0.4);
  \fill[brown!55](axis cs:-0.125,-0.125)rectangle(axis cs:0.125,0.125);
  \fill[brown!55](axis cs:0.4,0.4)rectangle(axis cs:0.6,0.6);
  \fill[brown!55](axis cs:0.325,-0.375)rectangle(axis cs:0.475,-0.225);
  \fill[brown!55](axis cs:-0.375,0.325)rectangle(axis cs:-0.225,0.475);
  \addplot[blue,thick,mark=*,mark size=1pt]coordinates{
  (-1.00,-1.00)(-0.85,-0.85)(-0.70,-0.70)(-0.67,-0.62)(-0.79,-0.53)(-0.92,-0.44)
  (-0.91,-0.32)(-0.80,-0.18)(-0.63,-0.05)(-0.45,0.07)(-0.27,0.18)(-0.08,0.29)
  (0.09,0.41)(0.27,0.53)(0.30,0.66)(0.29,0.84)(0.34,0.95)(0.51,1.01)(0.71,1.03)(0.94,1.03)};
  \addplot[only marks,mark=triangle*,green!50!black,mark size=4pt]coordinates{(-1,-1)};
  \addplot[only marks,mark=square*,red,mark size=4pt]coordinates{(1,1)};
\end{axis}
\end{tikzpicture}
\caption{CT-03: trajetória de $(-1,-1)$ a $(1,1)$, livre de colisões. Retângulos: obstáculos.}
\label{fig:traj}
\end{figure}

\begin{figure}[H]\centering
\begin{tikzpicture}
\begin{axis}[width=0.8\textwidth,height=5.6cm,xlabel={tempo $t$ (s)},ylabel={distância $d$ (m)},
  xmin=0,xmax=20,ymin=0,ymax=3,grid=both,grid style={gray!20},title={CT-03: convergência ao alvo}]
  \addplot[blue,thick,mark=*,mark size=1pt]coordinates{
  (0,2.83)(1,2.62)(2,2.40)(3,2.33)(4,2.36)(5.1,2.40)(6.1,2.33)(7.1,2.15)(8.1,1.94)
  (9.1,1.73)(10.1,1.51)(11.1,1.30)(12.1,1.08)(13.1,0.86)(14.1,0.77)(15.1,0.73)
  (16.1,0.66)(17.2,0.49)(18.2,0.29)(19.2,0.09)(19.3,0.0)};
  \addplot[red,dashed]coordinates{(0,0.07)(20,0.07)};
\end{axis}
\end{tikzpicture}
\caption{CT-03: distância ao alvo \emph{versus} tempo (linha tracejada: tolerância de chegada).}
\label{fig:dist}
\end{figure}

\subsection{CT-06 --- calibração dos sensores infravermelhos}
\begin{table}[H]\centering
\caption{CT-06: leitura de IR --- previsto (folha de dados) \emph{versus} medido.}
\label{tab:calib}
\begin{tabularx}{\textwidth}{@{}X c c@{}}
\toprule
\textbf{Condição} & \textbf{Previsto} & \textbf{Medido} \\
\midrule
Espaço livre ($>20$~cm) & $\approx 29$ & $\approx 100$--$120$ \\
Obstáculo frontal ($\approx 3$~cm) & $\approx 150$ & $\approx 250$--$300$ \\
\bottomrule
\end{tabularx}
\end{table}

% =====================================================================
\section{Registro de defeitos}
% =====================================================================
A Tabela~\ref{tab:def} consolida os defeitos detectados durante a campanha; todos foram corrigidos e os casos afetados, reexecutados com sucesso.

\begin{table}[H]\centering
\caption{Registro de defeitos.}
\label{tab:def}
\renewcommand{\arraystretch}{1.25}
\begin{tabularx}{\textwidth}{@{}l X c l@{}}
\toprule
\textbf{Id} & \textbf{Descrição} & \textbf{Sever.} & \textbf{Status} \\
\midrule
DEF-01 & Aprisionamento em mínimo local: robô preso na 1ª caixa (campo potencial puro), oscilando no lugar. & Alta & Corrigido \\
DEF-02 & Limiares de IR calibrados pela folha de dados ($\approx$29) em vez da base real do piso ($\approx$110): robô rodopiava na largada. & Alta & Corrigido \\
DEF-03 & Perda de telemetria por \emph{buffering} de \texttt{stdout}. & Baixa & Corrigido \\
\bottomrule
\end{tabularx}
\end{table}

\noindent\textbf{Evidência do DEF-01 (telemetria do aprisionamento):}
\begin{lstlisting}[style=c]
pos=(-0.54,-0.54) dist=1.90 err=+0.00 IR[L/F/R]=120/280/115 rodas=-3.0/+3.0
pos=(-0.54,-0.54) dist=1.90 err=-0.01 IR[L/F/R]=118/279/110 rodas=+3.0/-3.0
\end{lstlisting}
\noindent Posição congelada e rodas em oscilação $-3/+3 \to +3/-3$ caracterizam o mínimo local. \textbf{Correção:} introdução da máquina de estados com contorno comprometido (estilo Bug). \textbf{Correção do DEF-02:} recalibração dos limiares ($s_0{=}100$, \textsc{trig}${=}210$, \textsc{clear}${=}160$). \textbf{Correção do DEF-03:} gravação em arquivo com \texttt{fflush} por linha.

% =====================================================================
\section{Métricas}
% =====================================================================
\begin{itemize}[noitemsep,topsep=2pt]
  \item \textbf{Taxa de aprovação:} 6/6 casos = \textbf{100\%} (CT-04 aprovado após correção).
  \item \textbf{Defeitos:} 3 detectados, 3 corrigidos (100\% de tratamento).
  \item \textbf{Cobertura de requisitos essenciais:} RF-01, RF-02, RF-03, RF-06, RF-07, RNF-01 e RT-01 verificados.
\end{itemize}

% =====================================================================
\section{Conclusão}
% =====================================================================
O software de navegação foi \textbf{APROVADO} no cenário especificado: conduz o robô ao alvo de forma livre de colisões, escapa de mínimos locais e respeita o período do laço de tempo real. Registram-se como \textbf{ressalvas} a cobertura geométrica limitada, o uso de GPS em simulação (a ser substituído por odometria no robô físico) e a ausência de medição formal de WCET --- itens encaminhados como trabalhos futuros.

\end{document}
